\begin{table}[htbp]
\centering
\small
\begin{tabular}{@{}lccc@{}}
\toprule
&总体&小结石患者&大结石患者\\
\midrule
\begin{tabular}[c]{@{}l@{}}治疗 $a$：\\开放手术\end{tabular}
&$78\%$（273/350）&$\mathbf{93\%}$（81/87）&$\mathbf{73\%}$（192/263）\\
\midrule
\begin{tabular}[c]{@{}l@{}}治疗 $b$：\\经皮肾镜取石术\end{tabular}
&$\mathbf{83\%}$（289/350）&$87\%$（234/270）&$69\%$（55/80）\\
\bottomrule
\end{tabular}
\caption{Simpson 悖论的经典例子。表中给出两种肾结石治疗方法的成功率［Bottou et al., 2013；Charig et al., 1986，表 I、表 II］。虽然治疗 $b$ 的总体成功率看起来更高（每列中的粗体数值为该列最大值），但无论是小结石患者还是大结石患者，治疗 $b$ 都不如治疗 $a$（见例 6.37 与第 9.2 节）。}
\end{table}
